\documentclass[11pt,oneside]{article}
\newcommand{\TopicTitle}{Apuntes del Tema 8: Descenso de Gradiente, Métodos Acelerados y Adaptativos}
\usepackage[a4paper,margin=2.6cm,headheight=15pt]{geometry}
\usepackage{fontspec}
\usepackage{polyglossia}
\setdefaultlanguage{spanish}
\setotherlanguage{english}
\setmainfont{Latin Modern Roman}
\setsansfont{Latin Modern Sans}
\setmonofont{Latin Modern Mono}[Scale=MatchLowercase]
\usepackage{amsmath,amssymb,mathtools}
\usepackage{graphicx}
\usepackage{xcolor}
\usepackage{booktabs}
\usepackage{tabularx}
\usepackage{array}
\usepackage{longtable}
\usepackage{enumitem}
\usepackage{listings}
\usepackage{fancyhdr}
\usepackage{titlesec}
\usepackage{microtype}
\usepackage{hyperref}
\usepackage{bookmark}
\usepackage{caption}

\definecolor{ink}{HTML}{243447}
\definecolor{accent}{HTML}{176B87}
\definecolor{codebg}{HTML}{F3F5F7}
\definecolor{codegray}{HTML}{65727E}

\hypersetup{
  colorlinks=true,
  linkcolor=accent,
  urlcolor=accent,
  citecolor=accent,
  pdftitle={\TopicTitle},
  pdfauthor={Foundations of Machine Learning and Deep Learning}
}
\setcounter{secnumdepth}{3}
\setcounter{tocdepth}{3}
\setlist[itemize]{leftmargin=1.7em,itemsep=2pt,topsep=4pt}
\setlist[enumerate]{leftmargin=1.9em,itemsep=2pt,topsep=4pt}
\setlength{\parindent}{0pt}
\setlength{\parskip}{5pt plus 1pt minus 1pt}
\titleformat{\section}{\Large\sffamily\bfseries\color{ink}}{\thesection}{0.8em}{}
\titleformat{\subsection}{\large\sffamily\bfseries\color{accent}}{\thesubsection}{0.8em}{}
\titleformat{\subsubsection}{\normalsize\sffamily\bfseries\color{ink}}{\thesubsubsection}{0.8em}{}
\titlespacing*{\section}{0pt}{2.0em}{0.7em}
\titlespacing*{\subsection}{0pt}{1.4em}{0.4em}
\titlespacing*{\subsubsection}{0pt}{1.0em}{0.3em}
\lstset{
  backgroundcolor=\color{codebg},
  basicstyle=\ttfamily\small,
  breaklines=true,
  breakatwhitespace=false,
  columns=fullflexible,
  frame=single,
  framerule=0.3pt,
  rulecolor=\color{codegray},
  numbers=left,
  numbersep=8pt,
  numberstyle=\scriptsize\color{codegray},
  showstringspaces=false,
  keepspaces=true,
  tabsize=4,
  keywordstyle=\color{accent}\bfseries,
  commentstyle=\color{codegray}\itshape,
  stringstyle=\color{purple}
}
\captionsetup{font=small,labelfont=bf}
\pagestyle{fancy}
\fancyhf{}
\fancyhead[L]{\small\sffamily\color{codegray}Foundations of Machine Learning and Deep Learning}
\fancyhead[R]{\small\sffamily\color{codegray}\nouppercase{\rightmark}}
\fancyfoot[C]{\thepage}
\renewcommand{\headrulewidth}{0.4pt}
\renewcommand{\footrulewidth}{0pt}
\newcommand{\topiccover}[1]{%
  \begin{titlepage}
    \centering
    \vspace*{2.5cm}
    {\Large\sffamily\color{accent}APUNTES UNIVERSITARIOS\par}
    \vspace{1.5cm}
    {\Huge\sffamily\bfseries\color{ink}Foundations of Machine Learning\\[0.25em]
    and Deep Learning\par}
    \vspace{1.8cm}
    {\rule{0.72\textwidth}{1pt}\par}
    \vspace{1.2cm}
    {\LARGE\sffamily\bfseries #1\par}
    \vfill
    {\large\sffamily\color{codegray}\today\par}
  \end{titlepage}
}

\begin{document}
\topiccover{Apuntes del Tema 8: Descenso de Gradiente, Métodos Acelerados y Adaptativos}
\tableofcontents
\clearpage
\markboth{Apuntes del Tema 8: Descenso de Gradiente, Métodos Acelerados y Adaptativos}{Apuntes del Tema 8: Descenso de Gradiente, Métodos Acelerados y Adaptativos}
\section{Apuntes del Tema 8: Descenso de Gradiente, Métodos Acelerados y Adaptativos}

\subsection{1. Fundamentos de Optimización en Deep Learning}

\subsubsection{1.1 Optimización Matemática vs. Aprendizaje Profundo}

\begin{itemize}

\item \textbf{Optimización Pura:} Su objetivo principal es minimizar una función objetivo \(\mathcal{L}(w)\) sobre un conjunto de datos de entrenamiento para reducir el error de entrenamiento.

\end{itemize}

\begin{itemize}

\item \textbf{Deep Learning:} Su verdadero objetivo es reducir el \textbf{error de generalización} (desempeño en datos no vistos del conjunto de prueba), combinando la minimización del error con técnicas para mitigar el sobreajuste (\emph{overfitting}).

\end{itemize}

\subsubsection{1.2 Retos en la Superficie de Pérdida}

\begin{enumerate}

\item \textbf{Mínimos Locales:} Puntos donde el gradiente es cero (\(\nabla f(x) = 0\)) y la matriz Hessiana es definida positiva. En redes profundas con alta dimensionalidad, la mayoría de los puntos críticos no son mínimos locales, sino puntos de silla.

\end{enumerate}

\begin{enumerate}

\item **Puntos de Silla (\emph{Saddle Points}):** Puntos donde el gradiente se anula, pero la matriz Hessiana posee autovalores tanto positivos como negativos (la función sube en algunas direcciones y baja en otras).

\end{enumerate}

\begin{enumerate}

\item **Desvanecimiento del Gradiente (\emph{Vanishing Gradients}):** Regiones extremadamente llanas de la superficie de pérdida (por ejemplo, en las zonas de saturación de \(\tanh(x)\) o la sigmoide) donde \(\Vert{}\nabla f(x)\Vert{} \approx 0\), ralentizando drásticamente la actualización de los parámetros.

\end{enumerate}

\begin{lstlisting}

Superficie de Pérdida

  │        /\

  │       /  \   <- Máximo Local

  │  ____/    \____

  │ /              \  <- Punto de Silla (Hessiana con autovalores mixtos)

  │/                \____

  └─────────────────────── Mínimo Local / Global



\end{lstlisting}

\subsubsection{1.3 Convexidad y Series de Taylor}

\begin{itemize}

\item \textbf{Convexidad:} Un conjunto \(\mathcal{X}\) es convexo si el segmento que une cualquier par de puntos \(a, b \in \mathcal{X}\) está totalmente contenido en \(\mathcal{X}\). Una función \(f\) es convexa si:

\end{itemize}

\[\lambda f(x) + (1-\lambda)f(x') \ge f(\lambda x + (1-\lambda)x'), \quad \forall \lambda \in [0, 1]\]

\begin{itemize}

\item \emph{Propiedad clave:} En funciones estrictamente conexas, todo mínimo local es también un \textbf{mínimo global}. Aunque las funciones de pérdida en redes profundas son no convexas, exhiben comportamiento convexo en las vecindades inmediatas de los mínimos locales.

\end{itemize}

\begin{itemize}

\item \textbf{Expansión en Serie de Taylor:} Aproximación polinomial local de una función diferenciable \(f(x)\) alrededor de un punto \(x_0\):

\end{itemize}

\[f(x) \approx f(x_0) + \nabla f(x_0)^\top (x - x_0) + \frac{1}{2}(x - x_0)^\top \mathbf{H}(x - x_0)\]

donde \(\mathbf{H}\) es la matriz Hessiana de segundas derivadas parciales (\(\mathbf{H}_{ij} = \frac{\partial^2 f}{\partial x_i \partial x_j}\)).

\subsection{2. Algoritmos de Descenso de Gradiente}

\subsubsection{2.1 Descenso de Gradiente Unidimensional y Multivariante}

Partiendo de la expansión de Taylor de primer orden \(f(x + \epsilon) \approx f(x) + \epsilon \nabla f(x)\), para garantizar la máxima disminución del valor de \(f\), debemos movernos en la dirección opuesta al gradiente (\(\epsilon = -\eta \nabla f(x)\) con \(\eta > 0\)):

\[w_{t+1} = w_t - \eta \nabla f(w_t)\]

\begin{itemize}

\item \textbf{Impacto de la Tasa de Aprendizaje (\(\eta\)):}

\item Si \(\eta\) es demasiado pequeño: Convergencia extremadamente lenta.

\end{itemize}

\begin{itemize}

\item Si \(\eta\) es demasiado grande: Oscilaciones inestables e incapacidad para converger o divergencia.

\end{itemize}

\subsubsection{2.2 Variantes del Descenso de Gradiente según el Tamaño del Lote}

\begin{center}
\begin{tabularx}{\textwidth}{>{\raggedright\arraybackslash}X >{\raggedright\arraybackslash}X >{\raggedright\arraybackslash}X >{\raggedright\arraybackslash}X}
\toprule
Variante & Tamaño del Lote (\emph{Batch Size}) & Ventajas & Inconvenientes \\
\midrule
\textbf{Batch Gradient Descent} & Muestras totales \(N\) & Cálculo exacto del gradiente y convergencia suave. & Inviable en memoria RAM/GPU para datasets grandes; computacionalmente lento. \\
\bottomrule
\end{tabularx}
\end{center}

 |

| \textbf{Stochastic Gradient Descent (SGD)} | \(1\) muestra (\(\vert{}B\vert{} = 1\)) | Actualizaciones muy rápidas; la varianza del estimador ayuda a escapar de mínimos locales.

 | Gradientes con altísima varianza, trayectoria en "zig-zag" muy ruidosa.

 |

| \textbf{Mini-Batch SGD} | Lote pequeño (\(\vert{}B\vert{} \in [32, 512]\)) | \textbf{Estándar en la práctica:} Aprovecha la paralelizació en GPU/vectorización y suaviza el gradiente.

 | Requiere afinar la tasa de aprendizaje global \(\eta\).

 |

\subsubsection{2.3 Método de Newton (Segundo Orden)}

Aprovecha la curvatura local mediante la matriz Hessiana para calcular el paso óptimo:

\[\Delta x = -\mathbf{H}^{-1} \nabla f(x)\]

\begin{itemize}

\item \textbf{Ventaja:} Ajusta automáticamente la amplitud del paso sin depender fuertemente de \(\eta\).

\end{itemize}

\begin{itemize}

\item \textbf{Desventaja:} Computacionalmente prohibitivo en Deep Learning (\(O(d^3)\) para invertir \(\mathbf{H}\) con millones de parámetros) y sensible a direcciones donde la Hessiana es definida negativa.

\end{itemize}

\subsection{3. Optimizadores Acelerados y Adaptativos}

\begin{lstlisting}

Evolución de Optimizadores:

SGD ───► Momentum (Inercia v_t) ───► AdaGrad (Escalado por g²) ───► RMSProp (Promedio Móvil) ───► Adam (Momentum + RMSProp)



\end{lstlisting}

\subsubsection{3.1 Método del Momento (\emph{Momentum})}

Introduce una variable de velocidad \(v_t\) que acumula los gradientes pasados como un promedio móvil ponderado exponencialmente (física de una bola rodando cuesta abajo):

\[v_t \leftarrow \beta v_{t-1} + g_t\]

\[w_t \leftarrow w_{t-1} - \eta v_t\]

\begin{itemize}

\item \textbf{Propiedades:} Amortigua las oscilaciones en direcciones de alta varianza y acelera el avance a lo largo de barrancos de la superficie de pérdida. Típicamente \(\beta \approx 0.9\).

\end{itemize}

\subsubsection{3.2 AdaGrad (\emph{Adaptive Gradient Algorithm})}

Adapta la tasa de aprendizaje de cada parámetro individualmente basándose en la suma acumulada de sus gradientes al cuadrado \(s_t\):

\[s_t \leftarrow s_{t-1} + g_t \odot g_t\]

\[w_t \leftarrow w_{t-1} - \frac{\eta}{\sqrt{s_t + \epsilon}} \odot g_t\]

\begin{itemize}

\item \textbf{Propiedades:} Parámetros con actualizaciones frecuentes o grandes reciben tasas de aprendizaje efectivas más pequeñas, mientras que parámetros con actualizaciones raras o pequeñas reciben pasos más grandes.

\end{itemize}

\begin{itemize}

\item \textbf{Inconveniente:} \(s_t\) crece de forma monótona, haciendo que la tasa de aprendizaje efectiva tienda a cero y detenga el aprendizaje prematuramente.

\end{itemize}

\subsubsection{3.3 RMSProp y AdaDelta}

Resuelven la acumulación desmedida de AdaGrad sustituyendo la suma por un promedio móvil con decaimiento exponencial (\emph{leaky average}) con parámetro \(\gamma \in (0, 1)\):

\[s_t \leftarrow \gamma s_{t-1} + (1 - \gamma) g_t \odot g_t\]

\[w_t \leftarrow w_{t-1} - \frac{\eta}{\sqrt{s_t + \epsilon}} \odot g_t\]

\subsubsection{3.4 Adam (\emph{Adaptive Moment Estimation})}

El optimizador más popular y robusto en la práctica. Combina los beneficios del \textbf{Momento} (estimación del primer momento / media \(v_t\)) y de \textbf{RMSProp} (estimación del segundo momento / varianza \(s_t\)), añadiendo una \textbf{corrección de sesgo} para evitar la deriva hacia cero durante las primeras iteraciones:

\begin{enumerate}

\item \textbf{Primer Momento (Momento):}

\end{enumerate}

\[v_t \leftarrow \beta_1 v_{t-1} + (1 - \beta_1) g_t\]

\begin{enumerate}

\item \textbf{Segundo Momento (RMSProp):}

\end{enumerate}

\[s_t \leftarrow \beta_2 s_{t-1} + (1 - \beta_2) (g_t \odot g_t)\]

\begin{enumerate}

\item **Corrección de Sesgo (\emph{Bias Correction}):**

\end{enumerate}

\[\hat{v}_t = \frac{v_t}{1 - \beta_1^t}, \quad \hat{s}_t = \frac{s_t}{1 - \beta_2^t}\]

\begin{enumerate}

\item \textbf{Actualización de Parámetros:}

\end{enumerate}

\[w_t \leftarrow w_{t-1} - \frac{\eta}{\sqrt{\hat{s}_t} + \epsilon} \odot \hat{v}_t\]

\begin{quote}Hiperparámetros por defecto estándar: \(\beta_1 = 0.9\), \(\beta_2 = 0.999\), \(\epsilon = 10^{-8}\).\end{quote}

\begin{quote}\end{quote}

\begin{quote}\end{quote}

\subsection{4. Polítas de Ajuste (\emph{Learning Rate Schedulers}) y Warmup}

\subsubsection{4.1 Planificadores de Tasa de Aprendizaje (\emph{Schedulers})}

Reducen la tasa de aprendizaje \(\eta\) a medida que avanza el entrenamiento para evitar sobrepasar el mínimo local y permitir un ajuste fino:

\begin{itemize}

\item \textbf{Polynomial Decay:} \(\eta_t = \eta_0 \left(1 - \frac{t}{T}\right)^p\).

\end{itemize}

\begin{itemize}

\item \textbf{Step / Factor Decay:} \(\eta_{t+1} = \eta_t \cdot \alpha\) con \(\alpha \in (0, 1)\) cada \(k\) épocas.

\end{itemize}

\begin{itemize}

\item \textbf{Cosine Annealing:} Modula la tasa siguiendo una curva cosenoidal decreciente.

\end{itemize}

\begin{itemize}

\item \textbf{ReduceLROnPlateau:} Reduce \(\eta\) por un factor cuando una métrica de validación deja de mejorar durante varias épocas.

\end{itemize}

\subsubsection{4.2 Calentamiento (\emph{Warmup})}

Consiste en iniciar con una tasa de aprendizaje muy pequeña durante las primeras épocas e incrementarla linealmente hasta el valor \(\eta_0\) antes de aplicar cualquier estrategia de decaimiento. Previene la desestabilización inicial de los pesos cuando los gradientes son excesivamente volátiles.

\subsection{5. Práctica y Aplicación en Código (PyTorch)}

A continuación se incluye una implementación completa en PyTorch que compara el uso de \textbf{SGD con Momentum}, \textbf{Adam} y un \textbf{Scheduler con Warmup}:

\begin{lstlisting}[language=Python]

import torch

import torch.nn as nn

import torch.optim as optim

from torch.optim.lr_scheduler import LambdaLR, ReduceLROnPlateau



# --------------------------------------------------------------------------

# 1. Definición de una Red de Prueba

# --------------------------------------------------------------------------

class DummyNet(nn.Module):

    def __init__(self, input_dim=10, num_classes=2):

        super(DummyNet, self).__init__()

        self.fc = nn.Sequential(

            nn.Linear(input_dim, 32),

            nn.ReLU(),

            nn.Linear(32, num_classes)

        )

        

    def forward(self, x):

        return self.fc(x)



# --------------------------------------------------------------------------

# 2. Configuración de Optimizadores en PyTorch

# --------------------------------------------------------------------------

def setup_optimizers(model):

    # a) SGD con Momentum y Weight Decay

    optimizer_sgd = optim.SGD(

        model.parameters(), 

        lr=0.01, 

        momentum=0.9, 

        weight_decay=1e-4

    )

    

    # b) Adam (Adaptive Moment Estimation)

    optimizer_adam = optim.Adam(

        model.parameters(), 

        lr=0.001, 

        betas=(0.9, 0.999), 

        eps=1e-8

    )

    

    return optimizer_sgd, optimizer_adam



# --------------------------------------------------------------------------

# 3. Implementación de Warmup + Cosine Decay Scheduler

# --------------------------------------------------------------------------

def get_warmup_scheduler(optimizer, warmup_epochs=5, max_epochs=50):

    def lr_lambda(epoch):

        if epoch < warmup_epochs:

            # Incremento lineal durante el warmup

            return float(epoch + 1) / float(warmup_epochs)

        # Decaimiento pasados los epochs de warmup

        return 0.5 * (1.0 + torch.cos(torch.tensor((epoch - warmup_epochs) / (max_epochs - warmup_epochs) * 3.14159)))

    

    return LambdaLR(optimizer, lr_lambda=lr_lambda)



# --------------------------------------------------------------------------

# 4. Ejemplo de Bucle de Entrenamiento

# --------------------------------------------------------------------------

def train_demo():

    model = DummyNet()

    optimizer = optim.Adam(model.parameters(), lr=1e-3)

    scheduler = get_warmup_scheduler(optimizer, warmup_epochs=5, max_epochs=20)

    criterion = nn.CrossEntropyLoss()

    

    # Entrada y etiquetas de prueba

    inputs = torch.randn(16, 10)

    labels = torch.randint(0, 2, (16,))

    

    print("--- Demostración de Entrenamiento con Scheduler ---")

    for epoch in range(1, 21):

        optimizer.zero_grad()

        outputs = model(inputs)

        loss = criterion(outputs, labels)

        loss.backward()

        optimizer.step()

        

        # Obtener LR actual

        current_lr = optimizer.param_groups[0]['lr']

        print(f"Época [{epoch:2d}/20] | Loss: {loss.item():.4f} | LR: {current_lr:.6f}")

        

        # Paso del scheduler al final de la época

        scheduler.step()



if __name__ == "__main__":

    train_demo()



\end{lstlisting}

¿Hay algún punto o demostración de este bloque de optimizadores sobre el que quieras profundizar o tienes alguna duda específica?
\end{document}
